\begin{frame}{Статьи для анализа: генерация кода по тестам с LLM}

\textbf{Подход 2: LLM пишет код по готовым тестам человека}

\vspace{0.3cm}

\begin{itemize}
    \item \textbf{Mathews \& Nagappan, 2024 (ASE)}
    \begin{itemize}
        \item TDD-промптинг: тесты в промпте повышают процент решённых задач
        \item \textit{Проблема:} модели часто игнорируют часть тестов при большом их числе (эффект «потери середины»)
    \end{itemize}

    \vspace{0.2cm}

    \item \textbf{Cui, 2026 (arXiv — ICLR'25)}
    \begin{itemize}
        \item WebApp1K: бенчмарк TDD для веб-приложений, 19 моделей
        \item \textit{Проблема:} даже лучшие модели теряют инструкции на длинных промптах, производительность падает на 20-30\%
    \end{itemize}

    \vspace{0.2cm}

    \item \textbf{Piya \& Sullivan, 2024 (LLM4Code)}
    \begin{itemize}
        \item LLM4TDD: итеративная генерация кода по тестам (LeetCode)
        \item \textit{Проблема:} 11.5\% задач остаются нерешёнными; качество тестов, написанных пользователем, критично влияет на результат
    \end{itemize}
\end{itemize}

\end{frame}